\documentclass[../../../include/open-logic-section]{subfiles}

\begin{document}

\olfileid{his}{set}{cantorplane}
\olsection{రేఖ, తలంపై కాంటర్ వాదన}

ష్రోడర్--బెర్న్‌స్టైన్ నిరూపణకు సంబంధించిన కొన్ని పరిస్థితులు \olref[his][set][pathology]{sec}లో చర్చించిన చరిత్రతో ముడిపడి ఉన్నాయి. 1877లో ఒక చతురస్రంలో ఉన్నన్ని బిందువులే దాని ఒక భుజంలోనూ ఉంటాయని కాంటర్ నిరూపించారని గుర్తుంచుకోండి. ఇక్కడ ఆయన (మొదట ప్రయత్నించిన) నిరూపణను చూపుతాం.

$\unitline$ అనేది ప్రమాణ రేఖఖండం, అంటే $[0,1]$ బిందువుల సమితి అనుకుందాం. $\unitsquare$ అనేది ప్రమాణ చతురస్రం, అంటే $\unitline
\times \unitline$ బిందువుల సమితి అనుకుందాం. ఈ సంకేతాలలో కాంటర్ $\cardeq{\unitline}{\unitsquare}$ అని నిరూపించారు. డెడెకిండ్‌కు ఆయన రాసిన లేఖలో ప్రధానంగా ఈ క్రింది వాదన ఉంది.

\begin{thm}\ollabel{thm:cantorplane}
$\cardeq{\unitline}{\unitsquare}$
\end{thm}

\begin{proof}[నిరూపణ: మొదటి భాగం.]
$a, b \in \unitline$లను స్థిరపరచండి. వాటిని ద్వియాంశ సంకేతంలో రాయండి: $0$లు, $1$ల అనంత క్రమాలు $a_1$, $a_2$, \dots, అలాగే $b_1$, $b_2$, \dots కోసం ఇలా ఉంటుంది:
\begin{align*}
a &= 0.a_1a_2a_3a_4\dots\\
b &= 0.b_1b_2b_3b_4\dots
\intertext{ఇప్పుడు ఈ విధంగా ఇచ్చిన $f \colon \unitsquare \to \unitline$ ప్రమేయాన్ని చూడండి:}
f(a, b) & = 0.a_1b_1a_2b_2a_3b_3a_4b_4\dots
\end{align*}
ఈ నియమం నిజంగా ఒక ప్రమేయం కావడానికి ప్రతి సంఖ్యకు ఒకే ద్వియాంశ రూపాన్ని ఎంచుకుందాం: ఒకటి కన్నా చిన్న సంఖ్యలకు చివరికి అంతా ఒకట్లుగా మారని రూపాన్ని, ఒకటికి మాత్రం బిందువు తరువాత అన్నీ ఒకట్లున్న రూపాన్ని తీసుకుందాం. అప్పుడు $f$ ఒక \tetoken{అంతఃక్షేపణ}{!!a{injection}}: ఎందుకంటే $f(a, b) = f(c,d)$ అయితే అన్ని $n \in \Nat$లకు $a_n =
c_n$, $b_n = d_n$; అందువల్ల $a = c$, $b =
d$.
\sourcecorrection{OLTEHISSETCANP-001}{మూలం ద్వియాంశ రూపాలను ఏకైకంగా ఎంచుకోకుండానే ఈ ప్రమేయాన్ని, దాని అంతఃక్షేపణను ప్రకటించింది. ఏకైక ఎంపికను, ఒకటి అనే చివరి బిందువుకు అవసరమైన ప్రత్యేక రూపాన్ని స్పష్టం చేశాం.}
\end{proof}

దురదృష్టవశాత్తూ, డెడెకిండ్ కాంటర్‌కు చూపినట్లు, ఇంతటితో మొదటి ప్రశ్నకు సమాధానం రాదు. $0.\dot{1}\dot{0} =
0.1010101010\ldots$ను చూడండి. $f(a,b) = 0.\dot{1}\dot{0}$ కావాలంటే ఇక్కడ:
\begin{align*}
a&= 0.\dot{1}\dot{1} = 0.111111\ldots\\
b&= 0
\end{align*}
కానీ $a = 0.\dot{1}\dot{1} = 1$. కనుక ``$a$, $b$లను ద్వియాంశ సంకేతంలో రాయండి'' అనేటప్పుడు \emph{ఏ} రూపాన్ని వాడాలో ఎంచుకోవాలి. $f$ ఒక \emph{ప్రమేయం} కావాలి కనుక రెండు సాధ్య రూపాల్లో \emph{ఒక్కదాన్నే} వాడాలి. ఉదాహరణకు సరళంగా $a$ను ``$1.000\ldots$''గా రాస్తే, మూల వాదన ప్రకారం $f(a, b) = 0.\dot{1}\dot{0}$ అయ్యే $\tuple{a, b}$ జత లేదు. అయితే అప్పుడు పైన ఇచ్చిన సున్నా-పూర్ణభాగ రూపం ఒకటి వద్ద వర్తించదు; మనం ఎంచుకున్న చెల్లుబాటు అయ్యే నియమంలో ఈ ప్రత్యేక విలువకు పూర్వబింబం ఉంది.
\sourcecorrection{OLTEHISSETCANP-002}{మూలం ఇచ్చిన 0.1010... విలువ మనం స్పష్టం చేసిన మొత్తం-సమితి అంతఃక్షేపణలో f(1,0)గా వస్తుంది; అందువల్ల అది బింబం బయటకు ఉదాహరణ కాదు. నిజమైన వ్యతిరేక ఉదాహరణ కోసం, బేసి స్థానాల అంకెలు 0,1,1,1,... కాగా సరి స్థానాలవి 0,1,0,1,... అయ్యే ద్వియాంశ సంఖ్యను తీసుకోవచ్చు. దానికి ఏకైక ద్వియాంశ రూపం ఉంది; బేసి అంకెలు 1/2 యొక్క వాడని 0.0111... రూపం కావడంతో ఎంచుకున్న ప్రమేయానికి పూర్వబింబం ఉండదు.}

సారాంశంగా: కొన్ని పునరావృత్త ద్వియాంశ విస్తరణల సాధ్యత వల్ల, రూపం ఎంపికను స్థిరపరిచిన తరువాత కాంటర్ ప్రమేయం $f$ ఒక \tetoken{అంతఃక్షేపణ}{!!a{injection}}, కానీ \tetoken{అధిక్షేపణ}{!!a{surjection}} \emph{కాదు}. కనుక కాంటర్ చూపింది $\cardle{\unitsquare}{\unitline}$ మాత్రమే; $\cardeq{\unitsquare}{\unitline}$ \emph{కాదు}.

దాదాపు వెంటనే డెడెకిండ్‌కు తిరిగి రాస్తూ, కాంటర్ నిరూపణను ప్రధానంగా ఈ విధంగా పూర్తి చేయవచ్చని సూచించారు:

\begin{proof}[నిరూపణ: పూర్తిచేసిన భాగం.]
అంటే మనం $\cardle{\unitsquare}{\unitline}$ అని చూపాం. కానీ $\unitline$ నుంచి $\unitsquare$కు ఒక \tetoken{అంతఃక్షేపణ}{!!a{injection}} స్పష్టంగానే ఉంది: చతురస్రం ఒక భుజంపై రేఖను ఉంచితే చాలు. కాబట్టి $\cardle{\unitline}{\unitsquare}$, అలాగే $\cardle{\unitsquare}{\unitline}$. ష్రోడర్--బెర్న్‌స్టైన్ సిద్ధాంతం (\olref[sfr][siz][sb]{thm:schroder-bernstein}) ప్రకారం $\cardeq{\unitline}{\unitsquare}$.
\end{proof}

అయితే కాంటర్ ఆ చివరి పంక్తిని ఆ రూపంలో పూర్తి చేయలేకపోయేవారు; ఎందుకంటే అప్పటికి ష్రోడర్--బెర్న్‌స్టైన్ సిద్ధాంతం నిరూపించబడలేదు. తరువాత కాంటర్ దీన్ని సాధారణ పరికల్పనగా ప్రతిపాదించినా, సంతృప్తికరమైన నిరూపణ 1897 వరకు రాలేదు. (అందువల్ల 1877లో తరువాత కాంటర్ \olref{thm:cantorplane}కు ష్రోడర్--బెర్న్‌స్టైన్‌ను ఉపయోగించని మరో నిరూపణ ఇచ్చారు.)

\end{document}
